\begin{papercard}{lombard2002}
\cf{Stages} \stagetag{I} (also \stagetag{C} for pilot design and \stagetag{R} for the reporting checklist). A review plus a content analysis of 200 content-analysis articles (1994--1998) that shows how often intercoder reliability is not assessed or not reported, followed by ten numbered guidelines for assessing and reporting it. Krippendorff (2004) rebuts several of its statistical claims; read the two together.

\cf{Definitions} Intercoder reliability (more precisely \emph{intercoder agreement}, Tinsley \& Weiss) is the extent to which independent coders assign exactly the same code to a unit. It is distinct from correlational reliability, which only needs ratings to agree after standardisation \pc{lombard2002}{3}{intercoder agreement is needed in con}. Without it, content-analytic data and interpretations ``can never be considered valid''; it is necessary but not sufficient for validity, and it also permits dividing the coding work among coders \pc{lombard2002}{3}{establishing a high level of reliability also has the practical benefit of allowing the researcher to divide the coding work}. Latent content makes it \emph{more} important to show that the judgements are intersubjective.

\cf{Procedure assumed} Two or more coders categorise units from a representative set, independently and under the same conditions; a separate pilot test is used during training and a final test establishes the reliability of the full sample. Researchers acting as coders is questioned, as are overlapping unit assignments \pc{lombard2002}{4}{the coding decisions must be made independently under the same conditions}.

\cf{Indices reviewed} (1) \textbf{Percent agreement}: simple but ignores chance (two untrained coders agree on half of binary units); worst with few categories; can be inflated by adding rarely used categories; only for nominal data \pc{lombard2002}{4}{Without any instructions or training, without even knowing the property they are to identify, they will agree half of the time}. A reanalysis cited in the paper judged one-quarter to three-quarters of published observations unreliable once chance was corrected. (2) \textbf{Holsti's method}: identical to percent agreement when both coders code the same units; averaging it across variables ``hides'' weak variables \pc{lombard2002}{5}{a very poor practice which can hide variables with unacceptably low levels of reliability}. (3) \textbf{Scott's $\pi$}: uses pooled proportions, called ``too conservative'' because it presupposes identical coder distributions; Table 1 (81/9/9/1 of 100) gives $\pi=.05$. (4) \textbf{Cohen's $\kappa$}: uses multiplicative marginals; $\kappa=0.00$ for the same Table 1 \pc{lombard2002}{6}{kappa in the example illustrated in Table 1 would be 0.00}; maximum can be below 1 even at perfect agreement (Brennan \& Prediger); Fleiss for multiple coders; weighted $\kappa$ for partial credit. (5) \textbf{Krippendorff's $\alpha$}: any number of coders, any level of measurement, chance-corrected; main drawback is hand-calculation \pc{lombard2002}{6}{Krippendorff’s alpha index (1980) is attractive for several reasons}. Popping listed 39 nominal indices. No single ``best'' index exists; indices ignoring chance are ``too liberal'' while chance-corrected ones are ``too conservative''. Cronbach's $\alpha$ and $\chi^2$ are inappropriate \pc{lombard2002}{7}{There is also consensus that a few indices used are inappropriate measures of intercoder reliability}.

\cf{Acceptable levels} No established standard. Following Neuendorf: $\ge.90$ acceptable to all, $\ge.80$ acceptable in most situations, below that ``great disagreement''; $.70$ is often used for exploratory work, with more liberal criteria for the conservative indices \pc{lombard2002}{7}{coefficients of .90 or greater would be acceptable to all, .80 or greater would be acceptable in most situations}. Krippendorff (2004) disputes treating indices as ``liberal'' or ``conservative'' points on one scale.

\cf{Content analysis of content analyses} \emph{Sample}: all articles in \emph{Communication Abstracts} 1994--1998 with the keyword ``content analysis'' (census, $N=200$ after excluding 22 articles). \emph{Instrument}: informal independent coding of six articles, repeated and revised several times, then a formal pilot. \emph{Reliability sample}: the second and third authors both coded 128 articles (64\,\%) and later each coded half of the remaining 72; the final dataset used randomly chosen coder decisions for the doubly coded set \pc{lombard2002}{10}{To establish intercoder reliability, the second and third authors both coded 128 (64\%) of the articles}. Percent agreement, $\pi$, $\kappa$ (PRAM) and $\alpha$ (Krippendorff's Alpha 3.12) were computed per variable. \emph{Decision rule}: a variable was reliable if $\alpha\ge.70$, or, failing that, percent agreement $\ge.90$ \pc{lombard2002}{10}{be .70 or higher, or if this was not the case, percent agreement (a liberal index) be .90 or higher}.

\cf{Findings} Only 69\,\% ($n=137$) of the reports contained any reliability statement; they devoted a mean of 4.5 sentences to it, 6\,\% had a table, 45\,\% cited a reliability source, and the index was usually unspecified (when named: Holsti 15\,\%, Scott 10\,\%, percent agreement 9\,\%, Cohen 7\,\%, Krippendorff 3\,\%) \pc{lombard2002}{10}{Only 69\% of the research reports (n = 137) contained any report of intercoder reliability}. Missing information: size of the reliability sample (38\,\%), number of reliability coders (9\,\%), variable-level reliability (59\,\% reported only overall), training required (86\,\%), how discrepancies were resolved (74\,\%) \pc{lombard2002}{13}{the size of the reliability sample (missing in 38\% of the articles)}. Table 2 also shows how much coefficients differ for the same variable, e.g.\ for the rare medium ``Data from respondents'' (4 articles) percent agreement is 1.00 and $\kappa=1.00$ but $\pi=.66$ and $\alpha=.56$, and for ``Radio'' $\alpha=.67$ against 1.00 for the other three indices \pc{lombard2002}{11}{Radio 1.00 1.00 1.00 .67 3\% (5)}. The authors frame the problem as lack of guidelines and tools, not carelessness \pc{lombard2002}{13}{These results are not offered as an indictment of mass communication scholars or their work}.

\cf{Ten guidelines} \pc{lombard2002}{14}{Calculate and report intercoder reliability}:
(1) multiple coders and reported reliability; (2) choose one or more indices fitting the level of measurement, expected distributions and number of coders; if percent agreement is used, add a chance-corrected index and justify the choice; (3) obtain software; (4) set a minimum level in advance, $\ge.90$ nearly always acceptable, $\ge.80$ in most situations, $.70$ in exploratory work, with higher criteria for liberal indices and lower for conservative ones, ideally a two-index decision rule \pc{lombard2002}{14}{The preferred approach is to calculate and report two (or more) indices, establishing a decision rule}; (5) informal assessment during training on units outside the sample; (6) \emph{formal pilot} on a representative sample, rule of thumb \textbf{30 units}, independent coding, researcher preferably not a coder; if inadequate, retrain, refine the instrument, and only in extreme cases replace coders \pc{lombard2002}{15}{a good rule of thumb is 30 units}; (7) \emph{formal assessment on the full sample}: a random or otherwise justifiable sample of \textbf{at least 50 units or 10\,\%}, rarely more than 300, larger when the sample is large or expected reliability is low; pilot units may count only if pilot reliability was adequate \pc{lombard2002}{15}{should not be less than 50 units or 10\% of the full sample}; (8) choose and justify a procedure for incorporating the reliability coding into the final data (random selection, majority rule, expert tie-breaker, or discussion); (9) \emph{do not} use only percent agreement, correlation-based indices, $\chi^2$, an overall figure across variables, or overlapping reliability coding \pc{lombard2002}{16}{Use overall reliability across variables}.

\cf{Reporting checklist (guideline 10)} Reliability sample size, how it was created and why; its relation to the full sample (same, subset or separate); number of reliability coders and whether the researchers are among them; amount of coding by each coder; the index or indices and why; reliability per variable for each index; approximate training hours; how disagreements were resolved in the full sample; where to obtain the instrument, procedures and instructions \pc{lombard2002}{16}{The approximate amount of training (in hours) required to reach the reliability levels reported}.

\cf{Limits} Descriptive study of one field and period; the judgement of what counts as ``reported'' relied on the authors' own coding (themselves coders with moderate reliability on several variables, Table 2); the $.80/.70$ levels and the 30/50-unit sample sizes are rules of thumb without derivation; and the statistical claims about ``liberal'' and ``conservative'' indices and $\pi$'s assumptions are contested by Krippendorff (2004). Note also the guideline against overlapping coding, which Krippendorff (2004) argues is legitimate for $\alpha$.

\cf{Use for this thesis} The most compact checklist for the pilot and the report: record, \emph{before} coding, the pilot size (here the stage-C pilot should satisfy at least 30 units or justify fewer), the reliability-sample size on the main set ($\ge50$ units or 10\,\%), independence of the two coders, the second index beside percent agreement, per-code (not overall) values, training hours, and the disagreement-resolution rule. Treat its numeric thresholds as conventions from mass-communication research, not as app-store or privacy-specific standards, and justify any deviation. Use Krippendorff (2004) and Hallgren (2012) for the statistical choice, since this paper's own treatment of $\pi$, $\kappa$ and $\alpha$ is simplified.
\end{papercard}
